\section{AMC 8 2008}
        \begin{enumerate}
        	\item (\textit{AMC 8 2008}) Susan had 50 dollars to spend at the carnival. She spent 12 dollars on food and twice as much on rides. How many dollars did she have left to spend?


 $\textbf{(A)}\ 12 \qquad \textbf{(B)}\ 14 \qquad \textbf{(C)}\ 26 \qquad \textbf{(D)}\ 38 \qquad \textbf{(E)}\ 50$

 \textbf{Answer: }B



\textbf{Solution}

If Susan spent 12 dollars, then twice that much on rides, then she spent $12+12 \times 2=36$ dollars in total. We subtract $36$ from $50$ to get $\boxed{\textbf{(B)}\ 14}$



	\item (\textit{AMC 8 2008}) The ten-letter code $\text{BEST OF LUCK}$ represents the ten digits $0-9$, in order. What 4-digit number is represented by the code word $\text{CLUE}$?


 $\textbf{(A)}\ 8671 \qquad \textbf{(B)}\ 8672 \qquad \textbf{(C)}\ 9781 \qquad \textbf{(D)}\ 9782 \qquad \textbf{(E)}\ 9872$

 \textbf{Answer: }A



\textbf{Solution}

We can derive that $C=8$, $L=6$, $U=7$, and $E=1$. Therefore, the answer is $\boxed{\textbf{(A)}\ 8671}$~edited by OwenTang

 



\textbf{Solution 2}

We can easily see that $C=8$, which makes only $(A)$ $8761$ and $(B)$ $8762$ possible. Since the only difference between answer choices $A$ and $B$ is the last digit and that the last digit in the code word "CLUE" is $E$, we can just find that $E=1$ and that the answer is $\boxed{\textbf{(A)}\ 8671}$. ~solution by SuperVince1

 


	\item (\textit{AMC 8 2008}) If February is a month that contains Friday the $13^{\text{th}}$, what day of the week is February 1?


 $\textbf{(A)}\ \text{Sunday} \qquad \textbf{(B)}\ \text{Monday} \qquad \textbf{(C)}\ \text{Wednesday} \qquad \textbf{(D)}\ \text{Thursday}\qquad \textbf{(E)}\ \text{Saturday}$

 \textbf{Answer: }A



\textbf{Solution}

We can go backwards by days, but we can also backwards by weeks. If we go backwards by weeks, we see that February 6 is a Friday. If we now go backwards by days, February 1 is a $\boxed{\textbf{(A)}\ \text{Sunday}}$




\textbf{Solution 2}

Since the days of the week repeat every 7 days and $1-13\equiv2\pmod7$, the day of the week of the day 2 days after February 13 is the same as the day of the week of February 1. Since we know that February 13 is a Friday, the day of the week of the day 2 days after February 13 is a Sunday so we get that February 1 is a $\boxed{\textbf{(A)}\ \text{Sunday}}$.

 



\textbf{Solution 3}

We can go back 2 weeks and say that Friday is the -1st of February. Going forward 2 days, we get that February 1st is Sunday. So our answer is $\boxed{\textbf{(A)}\ \text{Sunday}}$.

 


	\item (\textit{AMC 8 2008}) In the figure, the outer equilateral triangle has area $16$, the inner equilateral triangle has area $1$, and the three trapezoids are congruent. What is the area of one of the trapezoids?

                    \begin{figure}[H]
                        \centering
                        \includegraphics[width=0.5\textwidth]{Problem/AMC_8_AMC_2008_4_Problem_Figure_1.png}
                    \end{figure}
                    
$\textbf{(A)}\ 3 \qquad \textbf{(B)}\ 4 \qquad  \textbf{(C)}\ 5 \qquad \textbf{(D)}\ 6 \qquad \textbf{(E)}\ 7$

 \textbf{Answer: }C



\textbf{Solution}

The area outside the small triangle but inside the large triangle is $16-1=15$. This is equally distributed between the three trapezoids. Each trapezoid has an area of $15/3 = \boxed{\textbf{(C)}\ 5}$.

 



\textbf{Solution 2}

You can roughly see that 1 trapezoid = 5 triangles = 5 Area since 1 triangle is 1 area so $\boxed{\textbf{(C)}\ 5}$



	\item (\textit{AMC 8 2008}) Barney Schwinn notices that the odometer on his bicycle reads $1441$, a palindrome, because it reads the same forward and backward. After riding $4$ more hours that day and $6$ the next, he notices that the odometer shows another palindrome, $1661$. What was his average speed in miles per hour?


 $\textbf{(A)}\ 15\qquad \textbf{(B)}\ 16\qquad \textbf{(C)}\ 18\qquad \textbf{(D)}\ 20\qquad \textbf{(E)}\ 22$

 \textbf{Answer: }E



\textbf{Solution}

Barney travels $1661-1441=220$ miles in $4+6=10$ hours for an average of $220/10=\boxed{\textbf{(E)}\ 22}$ miles per hour.

 


	\item (\textit{AMC 8 2008}) In the figure, what is the ratio of the area of the gray squares to the area of the white squares?

                    \begin{figure}[H]
                        \centering
                        \includegraphics[width=0.5\textwidth]{Problem/AMC_8_AMC_2008_6_Problem_Figure_1.png}
                    \end{figure}
                    
$\textbf{(A)}\ 3:10 \qquad\textbf{(B)}\ 3:8 \qquad\textbf{(C)}\ 3:7 \qquad\textbf{(D)}\ 3:5 \qquad\textbf{(E)}\ 1:1$

 \textbf{Answer: }D



\textbf{Solution}

Dividing the gray square into four smaller squares, there are $6$ gray tiles and $10$ white tiles, giving a ratio of $\boxed{\textbf{(D)}\ 3:5}$.

 


	\item (\textit{AMC 8 2008}) If $\frac{3}{5}=\frac{M}{45}=\frac{60}{N}$, what is $M+N$?


 $\textbf{(A)}\ 27\qquad \textbf{(B)}\ 29 \qquad \textbf{(C)}\ 45 \qquad \textbf{(D)}\ 105\qquad \textbf{(E)}\ 127$

 \textbf{Answer: }E



\textbf{Solution}

Separate into two equations $\frac35 = \frac{M}{45}$ and $\frac35 = \frac{60}{N}$ and solve for the unknowns. $M=27$ and $N=100$, therefore $M+N=\boxed{\textbf{(E)}\ 127}$.

 


	\item (\textit{AMC 8 2008}) Candy sales from the Boosters Club from January through April are shown. What were the average sales per month in dollars?

                    \begin{figure}[H]
                        \centering
                        \includegraphics[width=0.5\textwidth]{Problem/AMC_8_AMC_2008_8_Problem_Figure_1.png}
                    \end{figure}
                    
$\textbf{(A)}\ 60\qquad\textbf{(B)}\ 70\qquad\textbf{(C)}\ 75\qquad\textbf{(D)}\ 80\qquad\textbf{(E)}\ 85$

 \textbf{Answer: }D



\textbf{Solution}

There are a total of $100+60+40+120=320$ dollars of sales spread through $4$ months, for an average of $320/4 = \boxed{\textbf{(D)}\ 80}$.

 


	\item (\textit{AMC 8 2008}) In $2005$ Tycoon Tammy invested $100$ dollars for two years. During the first year
her investment suffered a $15\%$ loss, but during the second year the remaining
investment showed a $20\%$ gain. Over the two-year period, what was the change
in Tammy's investment?


 $\textbf{(A)}\  5\%\text{ loss}\qquad \textbf{(B)}\ 2\%\text{ loss}\qquad \textbf{(C)}\ 1\%\text{ gain}\qquad \textbf{(D)}\ 2\% \text{ gain} \qquad \textbf{(E)}\  5\%\text{ gain} \qquad$

 \textbf{Answer: }D



\textbf{Solution}

After the $15 \%$ loss, Tammy has $100 \cdot 0.85 = 85$ dollars. After the $20 \%$ gain, she has $85 \cdot 1.2 = 102$ dollars. This is an increase in $2$ dollars from her original $100$ dollars, a $\boxed{\textbf{(D)}\ 2 \%\ \text{gain}}$.

 


	\item (\textit{AMC 8 2008}) The average age of the $6$ people in Room A is $40$. The average age of the $4$ people in Room B is $25$. If the two groups are combined, what is the average age of all the people?


 $\textbf{(A)}\ 32.5 \qquad \textbf{(B)}\ 33 \qquad \textbf{(C)}\ 33.5 \qquad \textbf{(D)}\ 34\qquad \textbf{(E)}\ 35$

 \textbf{Answer: }D



\textbf{Solutions}

The total of all their ages over the number of people is

 
 \[\frac{6 \cdot 40 + 4\cdot 25}{6+4} = \frac{340}{10} = \boxed{\textbf{(D)}\ 34}.\]



	\item (\textit{AMC 8 2008}) Each of the $39$ students in the eighth grade at Lincoln Middle School has one dog or one cat or both a dog and a cat. Twenty students have a dog and $26$ students have a cat. How many students have both a dog and a cat?


 $\textbf{(A)}\ 7\qquad \textbf{(B)}\ 13\qquad \textbf{(C)}\ 19\qquad \textbf{(D)}\ 39\qquad \textbf{(E)}\ 46$

 \textbf{Answer: }A



\textbf{Solution 1}

The union of two sets is equal to the sum of each set minus their intersection. The number of students that have both a dog and a cat is $20+26-39 = \boxed{\textbf{(A)}\ 7}$.

 



\textbf{Solution 2 (Venn Diagram)}

We create a diagram:
                    \begin{figure}[H]
                        \centering
                        \includegraphics[width=0.5\textwidth]{Solution/AMC_8_AMC_2008_11_Solution_Figure_1.png}
                    \end{figure}
                    

 Let $x$ be the number of students with both a dog and a cat.

 Therefore, we have

 
 \[26+20-x = 39\]
 \[46-x = 39\]
 \[x = \boxed{\textbf{(A)} ~7}\].

 


	\item (\textit{AMC 8 2008}) A ball is dropped from a height of $3$ meters. On its first bounce it rises to a height of $2$ meters. It keeps falling and bouncing to $\frac{2}{3}$ of the height it reached in the previous bounce. On which bounce will it rise to a height less than $0.5$ meters?


 $\textbf{(A)}\  3 \qquad \textbf{(B)}\  4 \qquad \textbf{(C)}\ 5 \qquad \textbf{(D)}\ 6 \qquad \textbf{(E)}\ 7$

 \textbf{Answer: }C



\textbf{Solution}

Each bounce is $2/3$ times the height of the previous bounce. The first bounce reaches $2$ meters, the second $4/3$, the third $8/9$, the fourth $16/27$, and the fifth $32/81$. Half of $81$ is $40.5$, so the ball does not reach the required height on bounce $\boxed{\textbf{(C)}\ 5}$.

 


	\item (\textit{AMC 8 2008}) Mr. Harman needs to know the combined weight in pounds of three boxes he wants to mail. However, the only available scale is not accurate for weights less than $100$ pounds or more than $150$ pounds. So the boxes are weighed in pairs in every possible way. The results are $122$, $125$ and $127$ pounds. What is the combined weight in pounds of the three boxes?


 $\textbf{(A)}\ 160\qquad \textbf{(B)}\ 170\qquad \textbf{(C)}\ 187\qquad \textbf{(D)}\ 195\qquad \textbf{(E)}\ 354$

 \textbf{Answer: }C



\textbf{Solution 1}

Each box is weighed twice during this, so the combined weight of the three boxes is half the weight of these separate measures:

 
 \[\frac{122+125+127}{2} = \frac{374}{2} = \boxed{\textbf{(C)}\ 187}.\]




\textbf{Solution 2}

Using variables $a$, $b$, and $c$ to denote the boxes, with $a \le b \le c$. It is obvious that $a+b=122$ and $a+c=125$, since these are the two smallest pairs. Subtracting the former equation from the latter results in $c-b=3$. Additionally, $c+b=127$. Solving for $c$ and $b$ gives $c=65$ and $b=62$, so we can thus find $a=60$. Solving $60+62+65=\boxed{\textbf{(C) }\ 187}$




\textbf{Solution 3}

We use variables $x$, $y$, and $z$ to denote the three boxes. We know that $x+y=122$, $x+z=125$, and $y+z=127$. Note that the order of the variables or equations doesn't matter as long as they are consistent. By adding all three equations, we obtain $2x+2y+2z=122+125+127=374$. Dividing be 2, we obtain the sum $x+y+z$ is $\boxed{\textbf{(C) }\ 187}$.

 


	\item (\textit{AMC 8 2008}) Three $\text{A's}$, three $\text{B's}$, and three  $\text{C's}$ are placed in the nine spaces so that each row and column contain one of each letter. If $\text{A}$ is placed in the upper left corner, how many arrangements are possible?

                    \begin{figure}[H]
                        \centering
                        \includegraphics[width=0.5\textwidth]{Problem/AMC_8_AMC_2008_14_Problem_Figure_1.png}
                    \end{figure}
                    
$\textbf{(A)}\ 2\qquad\textbf{(B)}\ 3\qquad\textbf{(C)}\ 4\qquad\textbf{(D)}\ 5\qquad\textbf{(E)}\ 6$

 \textbf{Answer: }C



\textbf{Solution}

There are $2$ ways to place the remaining $\text{As}$, $2$ ways to place the remaining $\text{Bs}$, and $1$ way to place the remaining $\text{Cs}$ for a total of $(2)(2)(1) = \boxed{\textbf{(C)}\ 4}$.

 


	\item (\textit{AMC 8 2008}) In Theresa's first $8$ basketball games, she scored $7, 4, 3, 6, 8, 3, 1$ and $5$ points. In her ninth game, she scored fewer than $10$ points and her points-per-game average for the nine games was an integer. Similarly in her tenth game, she scored fewer than $10$ points and her points-per-game average for the $10$ games was also an integer. What is the product of the number of points she scored in the ninth and tenth games?


 $\textbf{(A)}\ 35\qquad \textbf{(B)}\ 40\qquad \textbf{(C)}\ 48\qquad \textbf{(D)}\ 56\qquad \textbf{(E)}\ 72$

 \textbf{Answer: }B



\textbf{Solution}

The total number of points from the first $8$ games is $7+4+3+6+8+3+1+5=37$. We have to make this a multiple of $9$ by scoring less than $10$ points. The smallest multiple of $9$ that is greater than $37$ is $45$. $45-37=8$. Now we have to add a number to get a multiple of 10. The next multiple of $10$ is $50$, which means that we add $5$. Multiplying these together, you get $8\cdot5 =$ $\boxed{\textbf{(B)}\ 40}$.

 


	\item (\textit{AMC 8 2008}) A shape is created by joining seven unit cubes, as shown. What is the ratio of the volume in cubic units to the surface area in square units?


 
                    \begin{figure}[H]
                        \centering
                        \includegraphics[width=0.5\textwidth]{Problem/AMC_8_AMC_2008_16_Problem_Figure_1.png}
                    \end{figure}
                    
$\textbf{(A)} \:1 : 6 \qquad\textbf{ (B)}\: 7 : 36 \qquad\textbf{(C)}\: 1 : 5 \qquad\textbf{(D)}\: 7 : 30\qquad\textbf{ (E)}\: 6 : 25$

 \textbf{Answer: }D



\textbf{Solution}

The volume is of seven unit cubes which is $7$. The surface area is the same for each of the protruding cubes which is $5\cdot 6=30$. The ratio of the volume to the surface area is $\boxed{\textbf{(D)}\: 7 : 30}$.

 


	\item (\textit{AMC 8 2008}) Ms. Osborne asks each student in her class to draw a rectangle with integer side lengths and a perimeter of $50$ units. All of her students calculate the area of the rectangle they draw. What is the difference between the largest and smallest possible areas of the rectangles?


 $\textbf{(A)}\ 76\qquad \textbf{(B)}\ 120\qquad \textbf{(C)}\ 128\qquad \textbf{(D)}\ 132\qquad \textbf{(E)}\ 136$

 \textbf{Answer: }D



\textbf{Solution}

A rectangle's area is maximized when its length and width are equivalent, or the two side lengths are closest together in this case with integer lengths. This occurs with the sides $12 \times 13 = 156$. Likewise, the area is smallest when the side lengths have the greatest difference, which is $1 \times 24 = 24$. The difference in area is $156-24=\boxed{\textbf{(D)}\ 132}$.

 


 


	\item (\textit{AMC 8 2008}) Two circles that share the same center have radii $10$ meters and $20$ meters. An aardvark runs along the path shown, starting at $A$ and ending at $K$. How many meters does the aardvark run?

                    \begin{figure}[H]
                        \centering
                        \includegraphics[width=0.5\textwidth]{Problem/AMC_8_AMC_2008_18_Problem_Figure_1.png}
                    \end{figure}
                    
$\textbf{(A)}\ 10\pi+20\qquad\textbf{(B)}\ 10\pi+30\qquad\textbf{(C)}\ 10\pi+40\qquad\textbf{(D)}\ 20\pi+20\qquad \\ \textbf{(E)}\ 20\pi+40$

 \textbf{Answer: }E



\textbf{Solution}

We will deal with this part by part:Part 1: 1/4 circumference of big circle= $\frac{2\pi r}{4}=\frac{\pi r}{2}=\frac{20\pi}{2}=10\pi$Part 2: Big radius minus small radius= $20-10=10$Part 3: 1/4 circumference of small circle= $\frac{\pi r}{2}=\frac{10\pi}{2}=5\pi$Part 4: Diameter of small circle: $2*10=20$Part 5: Same as part 3: $5\pi$Part 6: Same as part 2: $10$Total: $10\pi + 10 + 5\pi + 20 + 5\pi + 10 = \boxed{E = 20\pi + 40}$



	\item (\textit{AMC 8 2008}) Eight points are spaced around at intervals of one unit around a $2 \times 2$ square, as shown. Two of the $8$ points are chosen at random. What is the probability that the two points are one unit apart?

                    \begin{figure}[H]
                        \centering
                        
\includegraphics[width=0.5\textwidth]{Problem/AMC_8_AMC_2008_19_Problem_Figure_1.png}
                    \end{figure}
                    
$\textbf{(A)}\ \frac{1}{4}\qquad\textbf{(B)}\ \frac{2}{7}\qquad\textbf{(C)}\ \frac{4}{11}\qquad\textbf{(D)}\ \frac{1}{2}\qquad\textbf{(E)}\ \frac{4}{7}$

 \textbf{Answer: }B



\textbf{Solution 1}

The two points are one unit apart at $8$ places around the edge of the square. There are $8 \choose 2$$= 28$ ways to choose two points. The probability is

 
 \[\frac{8}{28} = \boxed{\textbf{(B)}\ \frac27}\]




\textbf{Solution 2}

Arbitrarily pick a point in the grid. Clearly, we see two options for the other point to be placed, so the answer is $\boxed{\textbf{(B)}\ \frac27}$.

 


	\item (\textit{AMC 8 2008}) The students in Mr. Neatkin's class took a penmanship test. Two-thirds of the boys and $\frac{3}{4}$ of the girls passed the test, and an equal number of boys and girls passed the test. What is the minimum possible number of students in the class?


 $\textbf{(A)}\ 12\qquad \textbf{(B)}\ 17\qquad \textbf{(C)}\ 24\qquad \textbf{(D)}\ 27\qquad \textbf{(E)}\ 36$

 \textbf{Answer: }B



\textbf{Solution}

Let $b$ be the number of boys and $g$ be the number of girls.

 
 \[\frac23 b = \frac34 g \Rightarrow b = \frac98 g\]

 For $g$ and $b$ to be integers, $g$ must cancel out with the denominator, and the smallest possible value is $8$. This yields $9$ boys. The minimum number of students is $8+9=\boxed{\textbf{(B)}\ 17}$.

 



\textbf{Solution 2}

We know that $\frac23 B = \frac34 G$  or  $\frac69 B = \frac68 G$. So, the ratio of the number of boys to girls is $8:9$ (Not $9:8$ because the numbers $8$ and $9$ are in the denominators of the fractions). The smallest total number of students is $9 + 8 = \boxed{\textbf{(B)}\ 17}$. ~DY, edited by SuperVince1

 


	\item (\textit{AMC 8 2008}) Jerry cuts a wedge from a $6$-cm cylinder of bologna as shown by the dashed curve. Which answer choice is closest to the volume of his wedge in cubic centimeters?

                    \begin{figure}[H]
                        \centering
                        \includegraphics[width=0.5\textwidth]{Problem/AMC_8_AMC_2008_21_Problem_Figure_1.png}
                    \end{figure}
                    
$\textbf{(A)} 48 \qquad \textbf{(B)} 75 \qquad \textbf{(C)}151\qquad \textbf{(D)}192 \qquad \textbf{(E)}603$

 \textbf{Answer: }C



\textbf{Solution 1}

The slice is cutting the cylinder into two equal wedges with equal volume. The cylinder's volume is $\pi r^2 h = \pi (4^2)(6) = 96\pi$. The volume of the wedge is half this which is $48\pi \approx \boxed{\textbf{(C)}\ 151}$.

 



\textbf{Solution 2}

The volume of the whole cylinder is $96\pi$. We approximate $\pi$ to be $3$, meaning our answer will be slightly too small. $96*3=288$, and we know the wedge cuts the cylinder perfectly in half. $288/2 = 144$, so the closest answer choice is $\boxed{\textbf{(C)}\ 151}$.

 


	\item (\textit{AMC 8 2008}) For how many positive integer values of $n$ are both $\frac{n}{3}$ and $3n$ three-digit whole numbers?


 $\textbf{(A)}\ 12\qquad \textbf{(B)}\ 21\qquad \textbf{(C)}\ 27\qquad \textbf{(D)}\ 33\qquad \textbf{(E)}\ 34$

 \textbf{Answer: }A



\textbf{Solution 1}

Instead of finding $n$, we find $x=\frac{n}{3}$. We want $x$ and $9x$ to be three-digit whole numbers. The smallest three-digit whole number is $100$, so that is our minimum value for $x$, since if $x \in \mathbb{Z^+}$, then $9x \in \mathbb{Z^+}$. The largest three-digit whole number divisible by $9$ is $999$, so our maximum value for $x$ is $\frac{999}{9}=111$. There are $12$ whole numbers in the closed set $\left[100,111\right]$ , so the answer is $\boxed{\textbf{(A)}\ 12}$.

 - ColtsFan10

 



\textbf{Solution 2}

We can set the following inequalities up to satisfy the conditions given by the question,$100 \leq \frac{n}{3} \leq 999$,and $100 \leq 3n \leq 999$.Once we simplify these and combine the restrictions, we get the inequality, $300 \leq n \leq 333$.Now we have to find all multiples of 3 in this range for $\frac{n}{3}$ to be an integer. We can compute this by setting $\frac{n} {3}=x$, where $x \in \mathbb{Z^+}$. Substituting $x$ for $n$ in the previous inequality, we get, $100 \leq x \leq 111$, and there are $111-100+1$ integers in this range giving us the answer, $\boxed{\textbf{(A)}\ 12}$.

 - kn07

 



\textbf{Solution 3}

We can create a list of the positive integers $n$ that fulfill the requirement of $\frac {n}{3}$ and $3n$ are three-digit whole numbers. The first number of this list must be $300$ since $\frac {300}{3} = 100$ is the smallest positive integer that satisfies this requirement. The last number of this list must be $333$ since $3 \cdot 333 = 999$ is the largest positive integer that satisfies this requirement. Since the problem requires $\frac {n}{3}$ and $3n$ must be whole numbers, the other numbers must be multiples of 3 (just like 300 and 333), so the list would look like this:
 \[300, 303, 306, . . . , 333\]To put this list in to a countable form we must put it in a form similar to $1,2,3, . . ., n$. So, we manipulate it as follows:
 \[300-300,303-300,306-300, . . .,333-300 \Rightarrow 0,3,6, . . ., 33\]

 \[\frac{0}{3}, \frac{3}{3}, \frac{6}{3}, . . ., \frac{33}{3} \Rightarrow 0,1,2, . . ., 11\]

 \[0+1,1+1,1+2, . . ., 11+1 \Rightarrow 1,2,3, . . ., 12\]

 Now we can tell that there are 12 positive integers which satisfies the two requirements, so the answer is $\boxed{\textbf{(A)}\ 12}$.

 



\textbf{Solution 4}

We have to find the range for all numbers first. We can first find the smallest digit in which $\frac {n}{3}$ is a 3 digit number that 3n is a 3 digit number which is $\frac {300}{3}$ = 100, then we find the biggest number in which $\frac {n}{3}$ is a 3 digit number and 3n is a 3 digit number which is $\frac {333}{3}$ = 111. Now we know that all integers from 100 to 111 are possible values and the amount of integer values in that range is $\boxed{\textbf{(A)}\ 12}$ (note I divided the numbers by 3 for convenience)

 


	\item (\textit{AMC 8 2008}) In square $ABCE$, $AF=2FE$ and $CD=2DE$. What is the ratio of the area of $\triangle BFD$ to the area of square $ABCE$?

                    \begin{figure}[H]
                        \centering
                        \includegraphics[width=0.5\textwidth]{Problem/AMC_8_AMC_2008_23_Problem_Figure_1.png}
                    \end{figure}
                    
$\textbf{(A)}\ \frac{1}{6}\qquad\textbf{(B)}\ \frac{2}{9}\qquad\textbf{(C)}\ \frac{5}{18}\qquad\textbf{(D)}\ \frac{1}{3}\qquad\textbf{(E)}\ \frac{7}{20}$

 \textbf{Answer: }C



\textbf{Solution 1}

The area of $\triangle BFD$ is the area of square $ABCE$ subtracted by the the area of the three triangles around it. Arbitrarily assign the side length of the square to be $6$.

 
                    \begin{figure}[H]
                        \centering
                        \includegraphics[width=0.5\textwidth]{Solution/AMC_8_AMC_2008_23_Solution_Figure_1.png}
                    \end{figure}
                    

 The ratio of the area of $\triangle BFD$ to the area of $ABCE$ is

 
 \[\frac{36-12-12-2}{36} = \frac{10}{36} = \boxed{\textbf{(C)}\ \frac{5}{18}}\]




\textbf{Solution 2}

Say that $\overline{FE}$ has length $x$, and that from there we can infer that $\overline{AF} = 2x$. We also know that $\overline{ED} = x$, and that $\overline{DC} = 2x$. The area of triangle $BFD$ is the square's area subtracted from the area of the excess triangles, which is simply these equations:
 \[9x^2 - (3x^2 + \dfrac{x}{2}^2 + 3x^2)\] 
 \[9x^2 - 6.5x^2\] 
 \[2.5x^2\]Thus, the area of the triangle is $2.5x^2$. We can now put the ratio of triangle $BFD$'s area to the area of the square $ABCE$ as a fraction. We have:
 \[\dfrac{2.5x^2}{9x^2}\]
 \[\dfrac{2.5\cancel{x^2}}{9\cancel{x^2}}\] 
 \[\dfrac{2.5}{9}\] 
 \[\dfrac{5}{18}\]Thus, our answer is $\boxed{\textbf{(C)}\ \frac{5}{18}}$



	\item (\textit{AMC 8 2008}) Ten tiles numbered $1$ through $10$ are turned face down. One tile is turned up at random, and a die is rolled. What is the probability that the product of the numbers on the tile and the die will be a square?


 $\textbf{(A)}\ \frac{1}{10}\qquad \textbf{(B)}\ \frac{1}{6}\qquad \textbf{(C)}\ \frac{11}{60}\qquad \textbf{(D)}\ \frac{1}{5}\qquad \textbf{(E)}\ \frac{7}{30}$

 \textbf{Answer: }C



\textbf{Solution}

The numbers can at most multiply to be $60$. The squares less than $60$ are $1,4,9,16,25,36,$ and $49$. The possible pairs are $(1,1),(1,4),(2,2),(4,1),(3,3),(9,1),(4,4),(8,2),(5,5),(6,6),$ and $(9,4)$. There are $11$ choices and $60$ possibilities giving a probability of $\boxed{\textbf{(C)}\ \frac{11}{60}}$.

 


	\item (\textit{AMC 8 2008}) Margie's winning art design is shown. The smallest circle has radius 2 inches, with each successive circle's radius increasing by 2 inches. Which of the following is closest to the percent of the design that is black?


 
                    \begin{figure}[H]
                        \centering
                        \includegraphics[width=0.5\textwidth]{Problem/AMC_8_AMC_2008_25_Problem_Figure_1.png}
                    \end{figure}
                    
$\textbf{(A)}\ 42\qquad \textbf{(B)}\ 44\qquad \textbf{(C)}\ 46\qquad \textbf{(D)}\ 47\qquad \textbf{(E)}\ 49\qquad$

 \textbf{Answer: }A



\textbf{Solution}

Let the smallest circle be 1, the second smallest circle be 2, the third smallest circle be 3, etc. 
 \[\begin{array}{c|cc} \text{circle \#} & \text{radius} & \text{area} \\ \hline 1 & 2 & 4\pi \\ 2 & 4 & 16\pi \\ 3 & 6 & 36\pi \\ 4 & 8 & 64\pi \\ 5 & 10 & 100\pi \\ 6 & 12 & 144\pi \end{array}\]

 The entire circle's area is $144\pi$. The area of the black regions is $(100-64)\pi + (36-16)\pi + 4\pi = 60\pi$. The percentage of the design that is black is $\frac{60\pi}{144\pi} = \frac{5}{12} \approx \boxed{\textbf{(A)}\ 42}$.

 


\end{enumerate}